% Architecture schematic for the MeshGraphNet / FluxGraphNet comparison.
% Requires: \usepackage{tikz}, \usetikzlibrary{arrows.meta,calc,positioning},
% and \usepackage{xcolor} (all are already loaded by this draft's preamble).
{%
\definecolor{archNavy}{HTML}{243B53}
\definecolor{archBlueTint}{HTML}{F1F5F9}
\definecolor{archLine}{HTML}{CBD5E1}
\definecolor{archText}{HTML}{334E68}
\definecolor{archOrange}{HTML}{C65D21}
\definecolor{archOrangeTint}{HTML}{FFF7ED}
\definecolor{archTeal}{HTML}{087F73}
\definecolor{archTealTint}{HTML}{ECFDF5}

\begin{figure}[htbp]
\centering
\begingroup
\linespread{1}\selectfont
\resizebox{0.78\linewidth}{!}{%
\begin{tikzpicture}[
    x=1cm,
    y=1cm,
    font=\footnotesize,
    sharedbox/.style={
        draw=archNavy,
        fill=archBlueTint,
        line width=0.8pt,
        rounded corners=3pt,
        align=center,
        minimum height=0.92cm,
        inner xsep=8pt,
        inner ysep=4pt
    },
    branchbox/.style={
        draw=archLine,
        fill=white,
        line width=0.65pt,
        rounded corners=3pt,
        align=center,
        minimum height=0.88cm,
        inner xsep=6pt,
        inner ysep=4pt
    },
    outputbox/.style={
        line width=0.9pt,
        rounded corners=3pt,
        align=center,
        minimum height=0.92cm,
        inner xsep=6pt,
        inner ysep=4pt
    },
    sharedarrow/.style={-{Latex[length=2.2mm,width=1.55mm]}, draw=archNavy, line width=1.05pt},
    mesharrow/.style={-{Latex[length=2.2mm,width=1.55mm]}, draw=archOrange, line width=1.05pt},
    fluxarrow/.style={-{Latex[length=2.2mm,width=1.55mm]}, draw=archTeal, line width=1.05pt}
]
    % Shared encode--process backbone.
    \node[
        fill=archNavy,
        text=white,
        rounded corners=2pt,
        font=\scriptsize\bfseries,
        inner xsep=7pt,
        inner ysep=3pt
    ] at (6.45,5.65) {SHARED ENCODE--PROCESS BACKBONE};

    \node[sharedbox, minimum width=3.20cm, text width=2.95cm, inner xsep=0pt]
        (input) at (2.25,4.60)
        {\textbf{Graph state}\\[-1pt]\color{archText}\scriptsize $u^n$ and edge geometry};
    \node[sharedbox, minimum width=3.05cm, text width=2.80cm, inner xsep=0pt]
        (encoder) at (5.825,4.60)
        {\textbf{Encoder}\\[-1pt]\color{archText}\scriptsize node + edge latents};
    \node[sharedbox, minimum width=4.45cm, text width=4.20cm, inner xsep=0pt]
        (processor) at (10.025,4.60)
        {\textbf{Message-passing processor}\\[-1pt]\color{archText}\scriptsize learned directional messages};

    \draw[sharedarrow] (input) -- (encoder);
    \draw[sharedarrow] (encoder) -- (processor);

    % Comparison panels
    \begin{scope}[yshift=-0.5cm]
    \path[
        draw=archOrange,
        fill=archOrangeTint,
        line width=0.8pt,
        rounded corners=5pt
    ] (0.15,-2.10) rectangle (6.25,3.15);
    \path[
        draw=archTeal,
        fill=archTealTint,
        line width=0.8pt,
        rounded corners=5pt
    ] (6.65,-2.10) rectangle (12.75,3.15);

    \node[
        fill=archOrange,
        text=white,
        rounded corners=2pt,
        font=\footnotesize\bfseries,
        minimum height=0.52cm,
        inner xsep=7pt,
        inner ysep=0pt
    ] (meshTitle) at (3.20,2.75) {MeshGraphNet};
    \node[
        fill=archTeal,
        text=white,
        rounded corners=2pt,
        font=\footnotesize\bfseries,
        minimum height=0.52cm,
        inner xsep=7pt,
        inner ysep=0pt
    ] (fluxTitle) at (9.70,2.75) {FluxGraphNet};

    \node[branchbox, minimum width=4.65cm, below=0.91cm of meshTitle]
        (nodeDecoder)
        {\textbf{Node decoder}\\[-1pt]\color{archText}\scriptsize decode each node embedding};
    \node[
        outputbox,
        draw=archOrange,
        fill=white,
        minimum width=4.65cm,
        below=0.91cm of nodeDecoder
    ] (nodeOutput)
        {\textbf{Direct node update}\\[-1pt]
         \color{archOrange}$\widehat{u}^{\,n+1}_i$ \quad for each $i\in\mathcal{V}$};

    \node[branchbox, minimum width=4.65cm, below=0.35cm of fluxTitle]
        (fluxDecoder)
        {\textbf{Interface-flux decoder}\\[-1pt]
         \color{archText}\scriptsize one oriented $\widehat{q}^{\,n}_{i+1/2}$ per interface};
    \node[branchbox, minimum width=4.65cm, minimum height=0.76cm,
        below=0.35cm of fluxDecoder]
        (fluxDifference)
        {\textbf{Conservative flux difference}};
    \node[
        outputbox,
        draw=archTeal,
        fill=white,
        minimum width=4.65cm,
        below=0.35cm of fluxDifference
    ] (fluxOutput)
        {\textbf{Conserved state}\\[-1pt]
         \color{archTeal}$\widehat{u}^{\,n+1}_i
         =u^n_i+\widehat{q}^{\,n}_{i-1/2}-\widehat{q}^{\,n}_{i+1/2}$};

    \end{scope}
    % A single shared representation feeds the two alternative decoders.
    %\coordinate (turn) at (10.025,3.93);
    %\coordinate (split) at (6.45,3.93);
    \coordinate (turn) at (10.025,3.60);
    \coordinate (split) at (6.45,3.60);
    \draw[sharedarrow, -] (processor.south) -- (turn) -- (split);
    \fill[archNavy] (split) circle (1.9pt);
    %\draw[mesharrow] (split) -- ++(-0.28,-0.48) -| (meshTitle.north);
    %\draw[fluxarrow] (split) -- ++(0.28,-0.48) -| (fluxTitle.north);
    \draw[mesharrow] (split) -- ++(-0.28,-0.25)
    -| (meshTitle.north);
    \draw[fluxarrow] (split) -- ++(0.28,-0.25)
        -| (fluxTitle.north);
    \draw[mesharrow] (meshTitle.south) -- (nodeDecoder.north);
    \draw[fluxarrow] (fluxTitle.south) -- (fluxDecoder.north);

    \draw[mesharrow] (nodeDecoder) -- (nodeOutput);
    \draw[fluxarrow] (fluxDecoder) -- (fluxDifference);
    \draw[fluxarrow] (fluxDifference) -- (fluxOutput);

\end{tikzpicture}%
}
\endgroup
\caption{Architectural comparison of MeshGraphNet and FluxGraphNet. Both models share the same encode--process backbone. MeshGraphNet decodes node-wise updates directly, whereas FluxGraphNet decodes one oriented transported amount per interface and applies a conservative flux-difference update, so each interface contribution enters neighboring cells with opposite signs. Detailed equations are given in Appendix \ref{app:model_details}.}
\label{fig:architecture_schematic}
\end{figure}
}%
